\documentclass{article}
\usepackage[T1]{fontenc}
\usepackage{lmodern}
\usepackage{graphicx}
\usepackage{amsmath,amssymb}
\usepackage[a4paper,margin=2cm]{geometry}
\usepackage[absolute,overlay]{textpos}
\setlength{\TPHorizModule}{1mm}
\setlength{\TPVertModule}{1mm}
\usepackage[indonesian]{babel}
\usepackage{hyperref}
\setlength{\emergencystretch}{2em}
% unit: Halaman 62

\begin{document}

% === Halaman 62 (llm) ===

\section*{Varian Vigenere Cipher}

Untuk mengatasi serangan dengan metode Kasiski, maka dibuat varian Vigenere Cipher sebagai berikut:

\begin{enumerate}
    \item \textcolor{red}{Full Vig\grave{e}n\grave{e}re cipher}
    \begin{itemize}
        \item Setiap baris di dalam tabel tidak menyatakan pergeseran huruf, tetapi merupakan permutasi huruf-huruf alfabet (lihat contoh table pada halaman berikut).
        \item Tabel tersebut harus dirahasiakan.
        \item Proses enkripsi dan dekripsi tetap sama seperti Vigenere standard:
    \end{itemize}
\end{enumerate}

\end{document}
